\subsection{Universal Constraint Systems (UCS)}\label{s:crypto_ucs}


To describe and motivate our framework we focus on two running examples. The first involves arithmetic modulo $2^{w}$, bitstring operations such as XOR and rotation, and arithmetic modulo an arbitrary integer (possibly non-prime, e.g.\ a power of two).  The second revolves around Hamming weight computations, lattice based operations, and a toy electronic voting application. %\kai{change end to, and has applications to areas such as voting}

We introduce some notation. Further background can be found in \cref{sec:preliminaries}. Given a ring or field $R$, we write $R^{<d}[X]$ to denote polynomials on the variable $X$ and of degree less than $d$. Note that $R^{<1}[X]=R$.   Given a bit size $w\ge 1$ we let $\Wordset{w}$ denote the set of integers $\{0,\ldots, 2^w-1\}$. We define the set of so-called \emph{bit-polynomials} of width $w$ as
%
\begin{equation}
\label{eq:def-bitpolyset}
  \BitPolyset{w} \;\coloneqq\; \Bigl\{\, \sum_{i=0}^{w-1} b_i\, X^i \;\Big|\; b_i \in \{0,1\} \,\Bigr\} \;\subseteq\; \ZZ[X].
\end{equation}
%
The set $\BitPolyset{w}$ is in bijective correspondence with $\Wordset{w}$ via the evaluation map $$\psi_2\colon\ZZ[X] \to \ZZ \qquad{\text{defined by}} \qquad \psi_2(u(X)) \coloneqq u(2),$$ which restricts to a bijection $\BitPolyset{w}\xrightarrow{\sim}\Wordset{w}$. Note that $\psi_2$ is the canonical surjective ring homomorphism $\ZZ[X] \twoheadrightarrow \ZZ[X]/(X-2)\cong \ZZ$. Here, $(X-2)$ denotes the ideal generated by $X-2$ in the ring $\ZZ[X]$. See \cref{sec:preliminaries} for background. We also let 
$$\phi_2\colon\ZZ[X] \to \FF_2[X]$$ denote reduction modulo $2$ coefficient-wise.

We emphasize that $\BitPolyset{w}$ is formally a set of \emph{integer} polynomials, and for the purposes of this paper this is not the same as $\FF_2^{<w}[X]$. It is important that commitments to polynomials in $\BitPolyset{w}$ are over $\ZZ[X]$, and cannot be taken modulo $2$. 

\subsubsection{A motivating example: Simultaneous bitstring operations and arithmetic mod $2^w$ and mod primes}

Our first example has  the following function $f$ at its core. Let $r\in [w-1]$ be a rotation parameter. Let $n\in \range{0}{2^w}$ be a an integer, possibly non-prime. Now define
\begin{align}
\label{eq:def-f-intro}
  f : \Wordset{w} &\to \Wordset{w},\nonumber\\
  x &\mapsto (x^2 \bmod n) \ \xor\ \ROTR^r(x),
\end{align}
where $\xor$ is the "bitwise exclusive or" (XOR) operation, and $\ROTR^r$ is right rotation by $r$ bits. Precisely,  for two integers $u,v\in\Wordset{w}$ with binary representations $u = \sum_{i=0}^{w-1} u_i 2^i$ and $v = \sum_{i=0}^{w-1} v_i 2^i$, we have $u \xor v = \sum_{i=0}^{w-1} (u_i + v_i \bmod 2) 2^i$, and $\ROTR^r(u) = \sum_{i=0}^{w-1} u_{(i+r) \bmod w} 2^i$. %\kai{should we explain $u_{i+r}$ for $i+r > w-1$ or is it well known} \albert{There's the $\mod w$ in the subscript, that should be clear enough :)} 
 We extend these operations to the set $\BitPolyset{w}$ by applying them coefficient wise. In the definition of $f$ (\cref{eq:def-f-intro}), the expression $x^2 \bmod n$ means that $x$ is understood as an integer in $\Wordset{w}$, squared, and then the canonical representative of $x^2$  modulo $2^w$ in $[0, n-1]$ is returned.

We view $f$ as a toy function roughly representative of computations involving some or all of the following: arithmetic modulo $2^w$, bitstring operations, and prime field operations. Concretely, we think of $f$ as a proxy for some of the subroutines in computations such as classic hashing, CPU and zkVM instructions, zkID authentication systems, TLS, etc. See  \cref{tab:computations}. %\kai{change Cf. to something else because it is right after etc.?}



Given a number of iterations $\numrows\geq 1$, our goal is to design a proof system for the  relation $\REL_f$ defined as
\begin{equation}
\REL_f = \left\{ (\inp;\wit)   \ \left| \  \begin{aligned} &\inp = (u_{1},u_\numrows) \in \Wordset{w}^2,\\ &\wit= (u_2,\ldots, u_{\numrows-1}) \in \Wordset{w}^{\numrows-2},\\   \ &u_{i+1} = f(u_i) \text{ for all } i \in [\numrows-1]  \end{aligned} \ \right.\right\}.
\end{equation}
%
Before diving into arithmetizing $\REL_f$, we make a series of observations about the operations used by the function $f$, which will motivate our arithmetization framework. 


%As a first step we describe $\REL_f$ by means of a series of polynomial constraints over suitable rings. Concretely, our arithmetization involves polynomial constraints over $\ZZ[X]$\footnote{Technically speaking, the ring is $\QQ[X]$ and not $\ZZ[X]$; see \cref{s:general_ucs_definition} for details.} and over the ring $\FF_2[X]$. Each constraint takes the form $Q(\ldots)\in \primeideal$, where $Q$ is a constraint polynomial and $\primeideal$ is an ideal of $\ZZ[X]$ or $\FF_2[X]$. When $\primeideal=(0)$ is the zero ideal, the constraint $Q(\ldots)\in \primeideal$ reduces to an algebraic equation $Q(\ldots)=0$. 



%We start by discussing each of the core operations involved in $f$: squaring modulo $n$, bitwise XOR, and rotation.  Throughout, we mostly work with elements from $\BitPolyset{w}$ rather than $\Wordset{w}$. 





\subparagraph{XOR and AND}  For any $u, v \in \BitPolyset{w}$, the following identity holds in $\ZZ[X]$:
\begin{equation}
\label{eq:xor-and-decomposition-intro}
  u+v 
  \;=\; \left(u\xor v\right) + 2\cdot \left(u\band v\right).
  \qquad\text{in }\ZZ[X],
\end{equation}
The converse also holds: if $u,v,z,a\in\BitPolyset{w}$ satisfy $u+v=z + 2\cdot a$ in $\ZZ[X]$, then $z=u\xor v$ and $a=u\band v$. Both directions follow from the fact that $b + b' = (b \xor b') + 2\,(b\band b')$  for any two $b,b'\in\{0,1\}$, applied coefficient-wise.

Hence, informally speaking, provided the witness elements belong to $\BitPolyset{w}$, the identity~\eqref{eq:xor-and-decomposition-intro} arithmetizes both XOR and AND over $\ZZ[X]$ simultaneously.

A further consequence of~\eqref{eq:xor-and-decomposition-intro} is that for any $u,v, z\in\BitPolyset{w}$,
\begin{equation}
\label{eq:xor-mod2-intro}
z = u \xor v \quad \Longleftrightarrow z = u + v  \mod 2,
\end{equation}
%
In summary, one can expose the XOR of two elements of $\BitPolyset{w}$ either via a single equality over $\ZZ[X]$ (which additionally yields the AND) or through addition modulo~$2$. We leverage both approaches in our arithmetization of SHA-256. 



\subparagraph{Rotation through ideal membership} For any $u\in\BitPolyset{w}$, the following ideal-membership relation holds in $\ZZ[X]$:
\begin{equation}
\label{eq:rotation-poly}
  \ROTR^r(u) - X^{w-r}\cdot  u \in (X^w-1).
\end{equation}
where $(X^w-1)$ denotes the ideal in $\ZZ[X]$ generated by the polynomial $X^w-1$. 
Conversely, for any two $u, v\in\BitPolyset{w}$, if $v - X^{w-r}\cdot u \in (X^w-1)$ then $v = \ROTR^r(u)$ (\cref{lem:rotation_ideal_membership}).
Hence rotation \emph{for elements in $\BitPolyset{w}$} can be enforced by the ideal-membership predicate $v - X^{w-r} u \in (X^w-1)$. 

At this point one may now wonder whether such a predicate can be efficiently handled in proof systems.  The answer is that this can be achieved by checking that certain appropriate multilinear extension, when evaluated at a random point, returns a value that belongs to the \kai{add "appropriate" since its not the same as the original ideal} \albert{It is! Isn't it?} ideal. This is analogous to the technique in classical proof systems where the MLE of the polynomial constraints composed with the MLE of the witnesses %\kai{is it MLE of the witness or something else? technically not MLE of a constraint} \albert{It's the MLE of the constraint } 
evaluates to zero at a randomly sampled point. Here, instead, evaluation to zero (i.e.\ membership to the zero ideal) is replaced by ideal membership. See \cref{lem:ideal_membership_mle_test}.

%\begin{remark}[Rotation over arbitrary rings]\label{rem:rot}
%\albert{could be removed} The rotation identity~\eqref{eq:rotation-poly} and the corresponding ideal-membership arithmetization hold verbatim for any polynomial $u\in R^{<w}[X]$ \albert{defined?} when $R$ is an arbitrary commutative ring, since the argument only uses the ring structure of $R[X]$.
%\end{remark}


\subparagraph{Mixing XOR and rotation}
%We now combine XOR and rotation. \kai{remove previous sentence to fit flow of other parts.} 
Given elements $u,v,z\in\BitPolyset{w}$, we want to arithmetize the constraint
\( z=u\xor \ROTR^{r}(v)\). To do so, observe that as long as $u,v, z$ belong to $\BitPolyset{w}$, the following equivalence holds:
%
\begin{equation}
\label{eq:xor-rot-ff2-intro}
  z=u\xor \ROTR^{r}(v) \quad \Longleftrightarrow \quad \phi_2(u) + \phi_2(v)\cdot X^{w-r} - \phi_2(z) \in (X^w-1) \qquad \text{in } \FF_2[X],
\end{equation}
%
%
% \conffalse
% \ifconf
% The equivalence follows from degree arguments and the XOR characterization~\eqref{eq:xor-mod2-intro}; see \cref{lem:sha256_sigma_char2} for a full proof.
% % \begin{comment}
% % \else
% % %
% % We sketch why the  converse direction of the equivalence holds; see \cref{lem:sha256_sigma_char2} for a full argument. Assume $u,v,z\in\BitPolyset{w}$  satisfy the right-hand side of \eqref{eq:xor-rot-ff2-intro} in $\FF_2[X]$. Write $\phi_2(v) X^{w-r} = \alpha +  (X^w-1)\beta$ with $\alpha\in \FF_2[X]$ of degree less than $w$, and $\beta\in\FF_2[X]$. Let $y\in \BitPolyset{w}\subseteq \ZZ[X]$ be such that $\phi_2(y)=\alpha$. Then the right-hand side of \cref{eq:xor-rot-ff2-intro} implies that $\phi_2(u) + \phi_2(y) - \phi_2(z)  \in (X^w-1)$ in $\FF_2[X]$. Since $u,z\in \BitPolyset{w}$, all elements in the right-hand side of this expression have degree less than $w$. Hence  $\phi_2(u) + \phi_2(y) - \phi_2(z) =0$ in $\FF_2[X]$. Since also $y\in \BitPolyset{w}$, it follows that $z= u\xor y$ in $\ZZ[X]$ (cf.\ \cref{eq:xor-rot-ff2-intro}). Finally, we have seen that $\phi_2(v) X^{w-r} - \phi_2(y) \in (X^w-1)$ in $\FF_2[X]$. By  \cref{eq:rotation-poly},  $y= \ROTR^r(v)$ in $\ZZ[X]$. This proves \cref{eq:xor-rot-ff2-intro}
% % \fi
% % \ifconf
% % \end{comment}
% \else
% \fi 
% %\kai{are we only describing the converse direction? maybe state it more directly if so.}

The equivalence in~\eqref{eq:xor-rot-ff2-intro} can be generalized to more complex expressions involving several XORs and rotations. It is particularly useful for arithmetizing some inner functions of SHA-256 and related hashes and ciphers.


\subparagraph{Linking bit-polynomials to integers, and squaring modulo an arbitrary integer $n$}  We would now like to, given $u,v\in\BitPolyset{w}$, express the constraint that the integer represented by $u$, i.e.\ $\psi_2(u)$, is the square of the integer represented by $v$, i.e.\ $\psi_2(v)$, modulo an arbitrary integer $n$. We can do so by observing the following: 
%
\begin{equation}\label{eq:mod-square-via-ideal}
\psi_2(v) = \psi_2(u)^2 \bmod n \quad \Longleftrightarrow \quad v-u^2 + n\cdot \mu \in (X-2) \qquad\text{in }\ZZ[X],
\end{equation}
for some $\mu \in \ZZ$.  To see why, assume there exists $\mu\in \ZZ$ such that  $v(X)-u(X)^2 + n\cdot \mu \in (X-2)$ in $\ZZ[X]$. Then $v(X)-u(X)^2 + n\cdot \mu$ vanishes at $X=2$, so $v(2)-u(2)^2 + n\cdot \mu = 0$ in $\ZZ$, which gives $\psi_2(v) = \psi_2(u)^2 \bmod n$. The forward direction is analogous.

A similar approach arithmetizes other algebraic modular operations.  An important underlying principle is that membership to the ideal $(X-2)$ links bit-representations (modeled as elements from $\BitPolyset{w}$) to the integer they represent, as long as elements are constrained to belong to $\BitPolyset{w}$ and $\Wordset{w}$. 



\begin{remark}[Lookup or typing constraints]\label{rem:correct_domain}
 In all examples above, the typing  of elements is crucial. I.e.\ most of the ideal membership constraints described above lose their intended effect if one does not pair them with \emph{lookup} or \emph{typing} constraints, e.g.\ the constraint \eqref{eq:mod-square-via-ideal} %\kai{missing word? [...] must be accompanied..} 
 must be accompanied of the lookup constraints $v,u\in \BitPolyset{w}$ and $\mu\in \ZZ$. Because of this, our arithmetizations consist mostly of algebraic constraints involving witness elements (these  may be zero equalities or ideal membership constraints), plus lookup or typing constraints.
\end{remark}


\subparagraph{Arithmetization of $\REL_f$} We now put together a full arithmetization for the relation $\REL_f$.
Recall that each application of the function $f$ computes
\[
  u_{t+1} \;=\; \bigl(u_{t}^2 \bmod n \bigr)\;\xor\; \ROTR^r\bigl(u_{t}\bigr).
\]
To arithmetize the full $(\numrows-1)$-fold iteration, we introduce witness vectors for each step 
and impose constraints that enforce the transition given by $f$.
The entire arithmetization uses three witness elements per step, two typed in $\BitPolyset{w}$ and one in $\ZZ$, and
two constraints  per transition. 

\medskip
\noindent\emph{\underline{Witness.}}
For each step $t\in[\numrows]$ we introduce three witness elements. We think of the witness as a $\numrows\times 3$ matrix, or \emph{trace} (even though our definition of UCS does not restrict to working with traces and AIR-like systems). 
 We denote the inverse of $\psi_2$ with $\hat{\cdot}$, so that for all $u\in \Wordset{w}$ we have $\hat{u}=\hat{u}(X)\in \BitPolyset{w}$ and $\hat{u}(2) = u$. 
The following table describes each witness element, its intended meaning, and its domain \albert{Change term ``domain'' to lookup or type, everywhere}.
\begin{center}
\begin{tabular}{lll}
\toprule
Entry & Intended meaning & Lookup constraint\\
\midrule
$\hat{u}_{t}$ & bit-polynomial representation of $u_{t}$ at step $t$ & $\in\BitPolyset{w}\subset\ZZ[X]$\\
$\hat{s}_{t}$ & bit-polynomial representation of $(u_{t})^2\bmod n$ & $\in\BitPolyset{w}\subset\ZZ[X]$\\
$\mu_{t}$ & the quotient such that $(u_{t})^2 = s_{t} + \mu_{t}\, n $ & $\in\ZZ$\\
\bottomrule
\end{tabular}
\end{center}


\medskip
\noindent\emph{\underline{Transition constraints.}}
For each step $t\in[\numrows-1]$ we impose two constraints.

\begin{enumerate}[leftmargin=*]
\item \textbf{Squaring modulo $n$.}
  \(
    \hat{s}_{t} - \bigl(\hat{u}_{t}\bigr)^2 + n\cdot \mu_{t} \;\in\; (X-2)
    \ \text{ in }  \ZZ[X].
  \)
  By \cref{eq:mod-square-via-ideal},
  this forces $s_{t}\equiv (u_{t})^2\pmod{n}$.

\item \textbf{XOR and rotation.}
  \(
    \phi_2\bigl(\hat{u}_{t+1}\bigr) - \phi_2\bigl(\hat{s}_{t}\bigr)
      - X^{w-r}\cdot\phi_2\bigl(\hat{u}_{t}\bigr)
      \;\in\; (X^w-1) \text{ in } \FF_2[X].
  \)
  By~\eqref{eq:xor-rot-ff2-intro}, this  enforces that $\hat{u}_{t+1} = \hat{s}_{t} \xor \ROTR^r(\hat{u}_{t})$.
\end{enumerate}
%
We omit boundary constraints for brevity. 

The reader may appreciate the proposed arithmetization as being both conceptually simple and small (in the sense of it requiring a small number of constraints given the relation sought to arithmetize). The question now is whether such arithmetization \emph{is useful in practice}, i.e.\ whether it can be efficiently handled by proof systems. One of the main contributions of this work is to show that the answer is yes. We outline why and how in the next sections of the present technical overview. 

% \begin{comment}

\medskip
\noindent\textbf{Boundary constraints.}
%The public input is the initial word $u_{1}\coloneqq\inp\in\Wordset{w}$.
%We enforce the linkage
%$\hat{u}_{1}(2)=\inp$
%via the boundary constraint $\hat{u}_{1} - \inp \in (X-2)$.
%The output is read off as $u_{\numrows} = \hat{u}_{\numrows}(2)$.
%Any satisfying witness then certifies $u_{\numrows} = f^{\numrows-1}(u_{1})$.

\cref{tab:relf_intro_summary} summarizes the UAIR${}^+$ trace layout and non-boundary constraint structure for the iterated computation $\REL_f$.

\begin{table}[h]
\centering
\begin{tabular}{rl|rl@{\;}c@{\;}l}
\toprule
\multicolumn{2}{c|}{\emph{Trace layout}} & \multicolumn{4}{c}{\emph{Constraints}} \\
\cmidrule(lr){1-2}\cmidrule(lr){3-6}
\#Cols & Type & \#Mon. & Ideal & & Ring\\
\midrule
$2$ & $\BitPolyset{w}$
& $3$ & $(X^w{-}1)$ & $\subseteq$ & $\FF_2[X]$ \\
$1$ & $\ZZ$
& $3$ & $(X{-}2)$ & $\subseteq$ & $\QQ[X]$ \\
\bottomrule
\end{tabular}
\caption{Summary of our arithmetization of $\REL_f$. The trace has  $\numrows$ rows (as many as iterations of $f$, plus one), and three columns, two typed in $\BitPolyset{w}$ and another in $\ZZ$.  The ``\#Mon.''\ column counts the total number of monomials across all constraints \cref{?}  sharing the same ideal and ring. All monomials are of degree $1$. Public-input, boundary constraints, and shifted row accessed (cf.\ \cref{?}) are omitted for simplicity.}
\label{tab:relf_intro_summary}
\end{table}
% \end{comment}


\subsubsection{Further examples}
\label{ss:further_examples}

The same ideal-membership approach extends to other settings; we highlight three (see \cref{sec:voting_example_appendix} for full worked-out examples).
\begin{itemize}[leftmargin=*]
  \item \textbf{Lattice-based operations.} Constraints in, for example, cyclotomic rings $R_q = \FF_q[X]/(\Phi(X))$ can be enforced through ideal-membership predicates over $\FF_q[X]$ where the ideal is $(\Phi(X))$. The typing of the witnesses can either be in some subset of  $\ZZ[X]$ or in $\FF_q[X]$,  depending on the application (see the discussion after \cref{def:ucs} for more details).
  \item \textbf{Hamming weight.} For any $u\in\BitPolyset{w}$, let $\wt(u)$ denote the number of nonzero coefficients of $u$. Then, given $c\in \ZZ$ we have that $\wt(u) = c \;\Longleftrightarrow\; u - c \in (X-1)$ in $\ZZ[X]$. This is because evaluation at $X=1$ sums coefficients (see \cref{lem:popcount}).
 \albert{Removed the voting example}
  %\item \textbf{Combination example.} See \cref{sec:voting_example_appendix} for an extended voting example combining Hamming-weight and lattice-based constraints.
\end{itemize}


\subsubsection{Universal Constraint Systems (UCS)}

We now introduce the main relation we will analyse. We call it the \emph{Universal Constraint System} (UCS) relation~$\ucs$. 


%t a high level, a UCS instance specifies a matrix of witness entries that are \emph{polynomials} (rather than field elements), together with constraints that assert ideal membership in various polynomial rings simultaneously. The key novelty over classical constraint systems (R1CS, CCS, AIR) is that a single witness vector~$\ff_0$ over~$\QQ[X]$ can participate in constraints over multiple rings---enabling, e.g., the same polynomial to drive both a modular-$2^w$ check (via the ideal~$(X-2)$) and a XOR operations bitstring rotation check (via~$(X^w-1)$).

A UCS instance is specified by an \emph{index} $\idx$ that fixes
a witness-vector length~$\witlen$,
a number of rows~$\numrows$,
a tuple of prime powers $\primetuple=(q_1,\ldots,q_{|\primetuple|})$,
a bit-size bound~$\bitbound$,
degree bounds~$\degbounds=(\degreebound_0,\ldots,\degreebound_{\numprimes})$,
and a set of constraints~$\constraints$ described below.

The witness consists of $\numprimes+1$ vectors $\ff_0, \ldots, \ff_{\numprimes}$, each of length~$\witlen$. The vectors are typed \albert{Ok let's use type for the IOPP enforced type and lookup as the type enforced by the PIOP} as
%
$$
\ff_0 \in \left(\mainpolyring{\QQ}{\bitbound}{\degreebound_0}\right)^{\witlen}, \qquad \ff_\primeindex \in \left(\FF_{q_\primeindex}^{<\degreebound_\primeindex}[X]\right)^{\witlen} \text{ for all } \primeindex\in[|\primetuple|],
$$
%
where, recall, $\mainpolyring{R}{\bitbound}{\degreebound}$ denotes the set of polynomials of degree $< \degreebound$ with coefficients in $R$, of bit-length less than $B$ (when $R$ is a field, we drop the bit bound and allow coefficients in all $R$). 
The entries of the witness vectors are thus \emph{polynomials themselves}. In particular, the multilinear extension (MLE) (cf.\ \cref{sec:multilinear-extensions}) of~$\ff_i$ is a multilinear polynomial whose coefficients lie in~$\QQ[X]$ or~$\FF_{q_\primeindex}[X]$---that is, a polynomial whose coefficients are polynomials. Regarding why $\ff_0$ is typed in $\mainpolyring{\QQ}{\bitbound}{\degreebound_0}$ rather than $\mainpolyring{\ZZ}{\bitbound}{\degreebound_0}$, see the technical observations after \cref{def:ucs}.

The public input is a vector $\yy\in(\mainpolyring{\QQ}{\bitbound}{\degreebound_0})^{\inplen}$.

Each constraint in~$\constraints$ is a tuple $(Q,\primeideal)$, where $Q$ is a polynomial with coefficients either in $\QQ$ or $\FF_{q_i}$ for some $i\in [\numprimes]$, and $\primeideal$ is an ideal of $\QQ[X]$ or $\FF_{q_i}[X]$, respectively. In the first case, $Q=Q(\rowvar,\yy,\ff_0)$ is a polynomial on $\ff_0$ and  $\yy$ and a row variable $\rowvar$ intended to range in $[\numrows]$. We call $(Q,\primeideal)$ a \emph{$\QQ[X]$-constraint}. In the later case, $Q=Q(\rowvar, \phi_{q_i}(\yy), \phi_{q_i}(\ff_0), \ff_i)$ is a polynomial on $\ff_i$, $\phi_{q_i}(\ff_0)$, $\phi_{q_i}(\yy)$ and $\rowvar$, where, recall, $\phi_{q_i}$ denotes (coefficient-wise) reduction modulo $q_i$.  We call this later $(Q,\primeideal)$ a \emph{$\FF_{q_i}$-constraint}. See \cref{rem:well-defined-projections} for a short comment on having $\phi_{q_i}(\ff_0)$, $\phi_{q_i}(\yy)$ be well-defined.

The relation~$\ucs$ consists of all triples $(\idx, \inp;\wit)$ such that, for every constraint $(Q,\primeideal)\in\constraints$ and every row $\rowindex\in[\numrows]$, the following holds: 
%
\begin{align*}
 &Q(\rowindex,\yy,\ff_0)\in\primeideal &&\text{ if } (Q,\primeideal) \text{ is a } \QQ[X]\text{-constraint},\\
  &Q(\rowindex,\phi_{q_i}(\yy),\phi_{q_i}(\ff_0),\ff_i)\in\primeideal &&\text{ if }  (Q,\primeideal) \text{ is a } \FF_{q_i}[X]\text{-constraint}, i\in [\numprimes].
\end{align*}
%
\begin{remark}[Shared witness principle]
The witness $\ff_0$ acts as a shared witness across different modular arithmetic.  This allowed in \cref{eq:def-f-intro} to have a single bit-polynomial column $\hat{u}\in\BitPolyset{w}\subset\ZZ[X]$ simultaneously participate in a squaring of integers modulo some integer, and XOR operations expressed as additions modulo $2$.
\end{remark}

The formal definition of UCS is given in \cref{def:ucs}. When restricted to a single finite field and the zero ideal, UCS subsumes standard constraint systems---R1CS, AIR, CCS, Plonkish, and lookup constraints---as special cases (cf.\ \cref{ss:ucs_r1cs,ss:ucs_lookup}).

After \cref{def:ucs} we give a detailed discussion of the design choices behind the definition of UCS, and further technical observations about the definition. 


\begin{remark}[Well-defined projections of $\ff_0$ and $\yy$ ]\label{rem:well-defined-projections} In the UCS definition, we also require that the reductions $\phi_{q_i}(\ff_0)$, $\phi_{q_i}(\yy)$ are well-defined. Namely, we ask that $\ff_0$ and $\yy$ have entries in $\mainpolyring{\local{q_i}}{\bitbound}{\degreebound_0}$, where, in short, $\local{q_i}$ is the subring of $\QQ$ formed by rationals whose denominator is not divisible by $q_i$  (see \cref{sec:local-rings}).
  
  
  In practice, this requirement is usually subsumed (i.e.\ implied) by the lookup or typing constraints on $\ff_0$ and $\yy$. For example, in our arithmetization of $\REL_f$, the witness $\ff_0$ contains bit-polynomial entries, which are polynomials with integer coefficients, which belong to $\local{q_i}$ for all prime or prime power $q_i$. 
\end{remark}

% \ifconf
% \subparagraph{Design choices.} The definition of UCS involves several deliberate design choices---in particular, why witness entries belong to polynomial rings $R^{<w}[X]$ rather than in $R^w$, and why we allow $\ZZ$/$\QQ$ witnesses rather than working entirely over a finite field (an integer computation can always be embedded into a large enough finite field via range constraints). In brief: the ring structure of $R[X]$ makes ideal-membership constraints natural, enables multiple ideals within the same relation, and supports random projection from large rings onto small fields---yielding significant compression. Working over $\ZZ$ further allows projecting onto different finite fields simultaneously (e.g.\ constraints modulo $2^{32}$, $2$, and a $256$-bit prime in our SHA-256 arithmetization) and avoids the overhead of large-field arithmetic throughout the proof. A detailed discussion of these design choices and further technical observations (lookup handling, typing in $\QQ$ vs.\ $\ZZ$, modular reduction of rationals) is given in \cref{s:general_ucs_definition}; see also the technical overview of the Zinc paper~\cite{Zinc}. %\albert{Referent remark about projection of lookups, which should maybe be written in the PIOP section.}
% \else

% \fi